\subsection{The full integral symplectic targets and their retained products}

\begin{statement}{E.56}{Proved}{Exact transvection targets at every modulus prime to six}
Let $J$ be the alternating chain matrix of E.52.7, put
$\Omega=J^{-1}$, and write
\[
 a=e_1,\qquad b=e_2,\qquad c=e_1+e_3,\qquad d=e_4,
 \qquad W=\langle a,b\rangle.
 \tag{E.56.1}
For $T=\mathbf Z/m$, where $(m,6)=1$, or
$T=R=\prod_{p\ge5}\mathbf Z_p$, define
\[
 \begin{split}
 T_v(t)&=I+t vv^{\mathsf t}J,\
 \mathcal P(T)&=\operatorname{Stab}_{\operatorname{Sp}_4(T,J)}(c),\
 \mathcal P^{\rm rt}(T)&=J\mathcal P(T)J^{-1}
                  \subset\operatorname{Sp}_4(T,\Omega).
 \end{split}
 \tag{E.56.2}
The four original column transvections
$T_a(1),T_b(1),T_{c-a}(1),T_d(1)$ generate
$\operatorname{Sp}_4(T,J)$; the first three generate
exactly $\mathcal P(T)$. For $T=R$, generation means
closed topological generation. The exact structure of
the latter group is
\[
 \begin{gathered}
 \mathcal P(T)=U(T)\rtimes\operatorname{SL}_2(T),\
 U(T)=\{n(w,z):w\in W(T),\ z\in T\},\
 n(w,z)=I+w c^{\mathsf t}J+c w^{\mathsf t}J
                         +z c c^{\mathsf t}J,\
 n(w,z)n(w',z')
     =n(w+w',z+z'+\langle w,w'\rangle).
 \end{gathered}
 \tag{E.56.3}
Here $\langle u,v\rangle=u^{\mathsf t}Jv$, and the
displayed $\operatorname{SL}_2$ acts on $W$ and fixes
$c,d$ individually.

In the literal right convention of E.52, the generators
are $M_i=T_{e_i}(1)^{\mathsf t}$. Their full image is
$\operatorname{Sp}_4(T,\Omega)$ and their first-three
image is exactly $\mathcal P^{\rm rt}(T)$. They fix
the row $c^{\mathsf t}$, or equivalently the column
$Jc=-e_4$. The precise source comparison is
\[
 J^{-1}M_iJ=T_{e_i}(-1),\qquad
 \operatorname{Ad}(J^{-1})\circ q^{\rm rt}_{T,n}
       =q^{\rm col}_{T,n}\circ\iota,
 \quad \iota(\sigma_i)=\sigma_i^{-1},\quad n=4,5.
 \tag{E.56.4}
The map $\iota$ is the actual involutive braid
automorphism given on all the displayed generators.

Retain the literal Ferrari maps to $T_L,T_{23}$ and
the exact quotient $S_L$ from E.38--E.43, and the
four- and five-strand Hurwitz maps to $A_{96}$ from
E.46. With their specified original generator maps,
the actual four-strand joint image is
\[
 \bigl(T_{23}\times_{S_L}T_L\bigr)
                 \times\mathcal P^{\rm rt}(T)\times A_{96},
 \tag{E.56.5}
\]
and the actual five-strand joint image is
\[
                  \operatorname{Sp}_4(T,\Omega)\times A_{96}.
 \tag{E.56.6}
In particular the actual joint image of the two
four-strand maps to $T_L$ and the new symplectic
target is $T_L\times\mathcal P^{\rm rt}(T)$.
All these statements retain their reductions and
the actual inverse-limit maps. The corrected words
of E.55 generate the same displayed joint targets.
\end{statement}

\begin{proof}
We give elementary matrix elimination and then determine
the common quotients through the specified braid maps.

\paragraph{The exact conventions and the Heisenberg subgroup.}
The integral matrix with columns $a,b,c,d$ has determinant
one and satisfies
\[
 (a\ b\ c\ d)^{\mathsf t}J(a\ b\ c\ d)
       =\begin{pmatrix}0&1\\-1&0\end{pmatrix}
                       \oplus
        \begin{pmatrix}0&1\\-1&0\end{pmatrix}.
 \tag{E.56.7}
\]
Thus $W$ and $\langle c,d\rangle$ are orthogonal
symplectic planes. Direct multiplication gives
$T_v(t)T_v(u)=T_v(t+u)$ and
$gT_v(t)g^{-1}=T_{gv}(t)$ for a symplectic $g$.
The first two transvections give the upper and lower
elementary matrices on $W$. For a unit $u$, put
\[
 w_W(u)=T_a(u)T_b(u^{-1})T_a(u).
 \tag{E.56.8}
\]
Its matrix on $W$ is
$\left(\begin{smallmatrix}0&u\\-u^{-1}&0\end{smallmatrix}\right)$,
and $w_W(u)w_W(-1)=\operatorname{diag}(u,u^{-1})$.

The half turn $h=w_W(1)^2$ sends $c-a$ to $c+a$
and fixes $c,d$. Orthogonality of $a,c$ gives the
exact polynomial identities
\[
 \begin{split}
 T_{c+a}(t)T_{c-a}(t)&=T_c(2t)T_a(2t),\
 T_{c+a}(t)T_a(-t)T_c(-t)
      &=I+t\bigl(a c^{\mathsf t}J+c a^{\mathsf t}J\bigr).
 \end{split}
 \tag{E.56.9}
\]
Consequently the first three original transvections
give $T_c(t)$ and $n(ta,0)$ for every parameter $t$.
Conjugation by $w_W(1)$ gives $n(tb,0)$ as well.
For finite $T$, every parameter is an integer power
of the corresponding unit-parameter transvection;
division by two in the first formula is legitimate
because two is a unit. Formula E.56.3 follows by
multiplication, and in particular
\[
 \begin{split}
 n(xa+yb,z)&=n(xa,0)n(yb,0)T_c(z-xy),\
 [n(xa,0),n(yb,0)]&=T_c(2xy).
 \end{split}
 \tag{E.56.10}
\]
These are actual words in the original three
generators, with every parameter and central term
retained.

To identify the whole stabilizer, a symplectic matrix
fixing $c$ induces a determinant-one matrix $A$ on
$c^\perp/Tc\simeq W$. Multiplying by the inverse
of the displayed lift of $A$ makes that induced
action the identity. The resulting matrix sends
\[
 c\longmapsto c,\qquad
 d\longmapsto d+w+zc,\qquad
 v\longmapsto v+\langle w,v\rangle c\quad(v\in W).
 \tag{E.56.11}
\]
The last formula is forced by pairing its image
with the image of $d$. Hence it is precisely
$n(w,z)$, uniquely. This proves E.56.3 through
the indicated matrix map.

Over a local ring $\mathbf Z/p^r$, every
$A=\left(\begin{smallmatrix}\alpha&\beta\\\gamma&\delta\end{smallmatrix}\right)$
of determinant one has a unit in its first column.
When $\alpha$ is a unit, the exact factorization is
\[
 A=E_{21}(\gamma/\alpha)
          \operatorname{diag}(\alpha,\alpha^{-1})
                                      E_{12}(\beta/\alpha).
 \tag{E.56.12}
\]
Otherwise left multiplication by $w_W(1)$ makes
the top entry of that column a unit. This proves
elementary generation of the required
$\operatorname{SL}_2$ and, with E.56.10, proves
that the first three original generators give
the full stabilizer over each local ring.

\paragraph{Eliminating a primitive vector.}
Adding $T_d(1)$ gives the same elementary
$\operatorname{SL}_2$ on $\langle c,d\rangle$.
We describe an actual word taking any primitive
vector
$v=v_W+zc+td$ to $c$ over $\mathbf Z/p^r$.
If neither $z$ nor $t$ is a unit, a coordinate
of $v_W$ is a unit. Choose $w=a$ or $b$ so that
$\langle w,v_W\rangle$ is a unit. Multiplication
by $n(w,0)$ changes the $c$-coordinate to
$z+\langle w,v_W\rangle$, which is then a unit.
If the $d$-coordinate is not yet a unit, the
plane rotation $w_{\langle c,d\rangle}(1)$
makes it one. Rename the resulting coordinates
$v_W,z,t$, with $t$ a unit.

Now apply, in this order,
\[
 n(-t^{-1}v_W,0),\quad T_c(-z/t),\quad
 \operatorname{diag}(t,t^{-1})\big|_{\langle c,d\rangle},\quad
 w_{\langle c,d\rangle}(1).
 \tag{E.56.13}
\]
The first map removes $v_W$ and leaves $z$ unchanged,
because $\langle -t^{-1}v_W,v_W\rangle=0$.
The next removes $z$; the last two send $td$ to
$d$ and then to $c$. Every factor has already been
expressed as a word in the four original transvections.

For an arbitrary symplectic $g$, apply this word
to $gc$. The resulting product fixes $c$, so it
belongs to the explicitly generated stabilizer.
This proves generation of the entire symplectic
group over each local ring. For a general $m$,
CRT parameters supported at a single prime-power
factor make each original transvection the identity
in all other factors. The local words therefore
give the whole product over $\mathbf Z/m$.
The stabilizer proof has the same property.
Compactness and the cofinal divisibility system
of moduli give the asserted closed images over $R$.

The transpose convention is also literal.
Since $J^{\mathsf t}=-J$, multiplication gives
$T_{e_i}(1)^{\mathsf t}=JT_{e_i}(-1)J^{-1}$.
The inverse-generator substitution preserves both
the Artin braid relations and the distant commutation
relations, and its square is the identity.
It is therefore the automorphism $\iota$ in E.56.4.
The equalities on every braid generator prove that
displayed morphism. They also prove both stated
fixed-vector descriptions, including $Jc=-e_4$.
The saved level-five and level-seven permutation
maps are compared with these transposes on every
nonzero vector in the certificate below.

\paragraph{Perfectness and the actual simple quotients.}
The groups just generated are perfect at every
finite level. On a symplectic plane, conjugation
by $\operatorname{diag}(2,1/2)$ multiplies the
upper elementary parameter by four; its commutator
with $E_{12}(t)$ is $E_{12}(3t)$.
The lower parameter is multiplied by $1/4$, and
the corresponding difference $-3/4$ is a unit.
Thus the plane $\operatorname{SL}_2$ is perfect.
Its commutators with $n(ta,0)$ give all those
elements, since the difference on $a$ is one.
Its difference on $b$ is $-1/2$, also a unit,
so the same holds for $n(tb,0)$.
Their commutators in E.56.10 give the whole
central subgroup. Hence $\mathcal P(T)$ is perfect.
For the full symplectic group, the primitive-vector
elimination conjugates each original transvection
to one in a symplectic plane; every generator is
therefore a commutator. This proves perfectness
of the full group too.

We include the finite-field simplicity argument
needed at arbitrary primes. The projective action
of $\operatorname{Sp}_4(\mathbf F_p,J)$ on lines
has kernel exactly $\{\pm I\}$: a linear operator
fixing all lines is scalar, and preservation of
the alternating form gives these two scalars.
The stabilizer of the line $\mathbf F_pc$ has
the three orbits of sizes
\[
                    1,\qquad p^2+p,\qquad p^3.
 \tag{E.56.14}
\]
For the perpendicular lines different from $c$,
the plane $\operatorname{SL}_2$ is transitive on
the nonzero $W$ component and $n(w,0)$ changes
the $c$ component arbitrarily. For a line not
perpendicular to $c$, scale its $d$ coordinate
to one; $n(w,z)$ then moves it to $d$.
This proves the orbit assertion using the actual
stabilizer already constructed.

A block containing the line of $c$ is a union
of these stabilizer orbits. The two intermediate
possible sizes are $p^2+p+1$ and $p^3+1$.
Neither divides the total number
$N=p^3+p^2+p+1$: one has
$N=p(p^2+p+1)+1$, and $p^3+1$ lies strictly
between $N/2$ and $N$. The faithful projective
action is consequently primitive.
The subgroup $\{T_c(t):t\in\mathbf F_p\}$ is
abelian and normal in the line stabilizer, since
replacing $c$ by $uc$ replaces $t$ by $u^2t$.
Its conjugates generate the group by the proved
transvection generation.

For completeness, the elementary normal-subgroup
argument now proves simplicity. A nontrivial normal
subgroup of a faithful primitive permutation group
is transitive. Modulo such a subgroup, every
conjugate of the displayed abelian subgroup has
the same image: write each group element as a
normal-subgroup element times an element of the
point stabilizer, which normalizes that subgroup.
The quotient is therefore generated by one
abelian image. Perfectness makes it trivial.
We have proved simplicity of
$\operatorname{PSp}_4(\mathbf F_p)$ for every
$p\ge5$. Simplicity of
$\operatorname{PSL}_2(\mathbf F_p)$ is proved
by the same normal-subgroup argument in E.53.

The kernel of reduction from a matrix group
modulo $p^r$ to its reduction modulo $p$ is
a $p$-group: successive kernels embed, by
$1+p^jA\mapsto A$, into additive matrix spaces
over $\mathbf F_p$.
The structure E.56.3 therefore shows that
\[
 \begin{split}
 \mathcal P(\mathbf Z/p^r)&\longrightarrow
                         \operatorname{PSL}_2(\mathbf F_p),\
 \operatorname{Sp}_4(\mathbf Z/p^r)&\longrightarrow
                         \operatorname{PSp}_4(\mathbf F_p)
 \end{split}
 \tag{E.56.15}
\]
have solvable kernels. These are the actual
maps induced on $c^\perp/Tc$, followed by
reduction and the scalar quotient, in the first
case, and the actual reduction and scalar
quotient in the second.

Every proper normal subgroup of either source
in E.56.15 is contained in the stated kernel.
Indeed its image in the simple target is
trivial or the whole target. In the latter
case the resulting quotient is an image of
the solvable kernel; it is also a quotient of
a perfect group and hence must be trivial.
Thus every nontrivial quotient retains the
specified simple quotient in E.56.15.
This proves the required assertion including
all central and higher congruence kernels.

\paragraph{The complete target product with $T_L$.}
Consider the actual subdirect image of the two
four-strand maps to $T_L$ and
$\mathcal P^{\rm rt}(\mathbf Z/m)$.
A nontrivial Goursat common quotient is perfect,
because the latter group is perfect. It has a
nonabelian simple quotient. The matrix target
is the product of its prime-power factors;
their images in a simple group commute and
are normal, so exactly one factor supplies
that simple quotient. E.56.15 identifies it
with the specified
$\operatorname{PSL}_2(\mathbf F_p)$ map, up
to an automorphism of that target.

For $p=5,7$, this would force that actual
projective matrix map to factor through
$T_L$. The complete literal products
$T_L\times\mathcal P^{\rm rt}(\mathbf F_p)$
proved in E.50 and E.51 exclude that factorization.
For $p\ge11$, the order of
$\operatorname{PSL}_2(\mathbf F_p)$ has a
factor $p$, whereas the retained exact order
of $T_L$ has prime divisors only $2,3,5,7$.
This excludes every possible common quotient.
Goursat therefore gives the whole actual
product $T_L\times\mathcal P^{\rm rt}(\mathbf Z/m)$.

The projection
$T_{23}\times_{S_L}T_L\to T_L$ has solvable
kernel. In a common quotient with the perfect
matrix target, quotienting by the image of
that kernel gives a common quotient with
$T_L$, just excluded. The original common
quotient would then be solvable and perfect,
and hence trivial. This proves the product
with the whole old four-strand target.

\paragraph{The alternating factor and the full inverse limit.}
The simple group $A_{96}$ is neither
$\operatorname{PSL}_2(\mathbf F_p)$ nor
$\operatorname{PSp}_4(\mathbf F_p)$ for any
$p\ge5$. The matrix group orders are
\[
 \frac{p(p^2-1)}2,\qquad
                 \frac{p^4(p^2-1)(p^4-1)}2.
 \tag{E.56.16}
\]
For the second formula, count a symplectic
basis: its successive choices number
$(p^4-1),p^3,(p^2-1),p$, and then divide by
the two scalar matrices. The first formula
is the same basis count in dimension two.
Equality with $96!/2$ would imply $p\le96$,
since $p$ divides that order. Both displayed
orders would then be smaller than $96^{10}/2$,
while $96!\ge77^{20}>96^{10}$.
This proves the asserted exclusions.

It follows from E.56.15 that no finite matrix
target in the statement has an $A_{96}$
quotient. A simple quotient of a product
comes from one of its factors, as above.
The old target has no such quotient either:
its solvable kernel is killed in a simple
nonabelian quotient, and the order of $T_L$
has no factor eleven. The actual four- and
five-strand Hurwitz maps are onto $A_{96}$
by E.46. Goursat now proves exactly
E.56.5 and E.56.6 at every finite modulus.

Every finite collection of congruence levels
is dominated by their least common multiple.
The closed actual images over $R$ therefore
surject onto every finite projection of the
displayed products. Compactness proves the
same product identities over $R$.

Finally E.55 retains the same Galois word in
all matrix and $T_L$ components and in the
Hurwitz component. Their corrected maps
are the displayed original joint maps
precomposed with one actual braid
automorphism, so their images are the same
products. The separate corrected $T_{23}$
map is onto and agrees in $S_L$ by E.55.
The maximum solvable quotient of $T_L$ is
the specified $S_L$; the new matrix and
alternating factors are perfect. Thus the
same solvable-kernel Goursat argument gives
the whole corrected product E.56.5.
All inverse-limit comparisons retain these
same maps. This proves the final assertion.
\end{proof}

The new certificate has the shared prefix
\texttt{bc\_symplectic\_integral\_targets}.
It verifies the polynomial identities in
E.56.9--E.56.10, the exact transpose and
mirror comparison, and all $12096$ saved
level-five and level-seven permutation entries.
It eliminates all $624+2400$ nonzero vectors
over those two fields by literal original
transvection words. It also constructs
independent symplectic bases and factors
$16$ matrices at each of the moduli
$5,7,25,49,35,175,245$, retaining a complete
word witness at each modulus. Its run takes
$3.076$ seconds. Every protected earlier
certificate is unchanged.
